%%
%% Copyright (c) 2018-2020 Weitian LI <wt@liwt.net>
%% CC BY 4.0 License
%%
%% Created: 2018-04-11
%%

% Chinese version
\documentclass[zh]{resume}

% === 强制一页:紧缩间距 ===
\usepackage{geometry}
\geometry{margin=0.9cm}           % 默认 1.5cm → 0.9cm
\linespread{1.1}               % 中文默认 1.25 → 1.1
\titlespacing{\section}{0em}{0.5ex}{0.5ex}  % 节前后间距减半
\setlist[itemize,1]{label=\faAngleRight, nosep, leftmargin=2em, topsep=1pt}
\setlength{\medskipamount}{4pt}

% File information shown at the footer of the last page
\fileinfo{%
  % \faCopyright{} 2025--2026, 阿东玩AI \hspace{0.5em}
  % \creativecommons{by}{4.0} \hspace{0.5em}
  % \githublink{dok4everak47}{LLM-Resume-Template} \hspace{0.5em}
  % \faEdit{} \today
}

\name{劲舟}{梁}

\tagline{求职意向: 数据科学方向（数据分析 / 特征工程 / 模型评估）}

\photo[shape=circular, position=right]{1.5cm}{../assets/luongchin.png}

\profile{
  \mobile{18176608865}
  \email{dok4ever123@126.com}
  \iconlink{\faGithub}{https://github.com/dok4everak47}{dok4everak47}
}

\begin{document}
\begin{tcolorbox}[colback=accentcolor!12, colframe=accentcolor!20,
  arc=0pt, boxrule=0.3pt, left=6pt, right=6pt, top=3pt, bottom=3pt,
  sharp corners]
\makeheader
\end{tcolorbox}

\vspace{0.3em}

%======================================================================
\sectionTitle{个人总结}{\faUser}
%======================================================================
具有数学与计算科学背景的数据方向候选人。硕士阶段打下扎实的数理统计与分析基础；实战中主导过多模态 AIGC 数据评估体系与 Agent 对话日志分析系统的搭建，能将复杂数据问题转化为可量化的分析方案，在数据质量管控、指标体系构建、特征提取和模型行为分析方面有深入的一线经验。


%======================================================================
\sectionTitle{专业技能}{\faWrench}
%======================================================================
\begin{itemize}
    \item \textbf{编程语言}: TypeScript, Python, Node.js, SQL
    \item \textbf{数据分析}: Pandas, NumPy, scikit-learn, 数据清洗, 特征工程, 统计建模
    \item \textbf{模型评估}: 多模态 AIGC 评估体系, Agent 行为对比评测, 标注质量控制, Benchmark 设计
    \item \textbf{工具链}: Git, Linux 命令行, Docker, Vim/Neovim, VSCode
\end{itemize}

%======================================================================
\sectionTitle{个人项目}{\faRocket}
%======================================================================
\begin{itemize}
    \item[\faCode] \textbf{Claude Code from Scratch — 大型 Agent 系统架构实现}
    | TypeScript, Node.js \\
    GitHub: \href{https://github.com/dok4everak47/claude-code-from-scratch}{\texttt{dok4everak47/claude-code-from-scratch}}

    从零复现 Coding Agent 核心架构（4300+ 行），实现 LLM 调用→工具执行→结果观察→循环调度的完整闭环，设计分层上下文压缩与多类型记忆系统。体现了复杂系统的工程设计与数据流架构能力。

    \vspace{1ex}
    {\color{accentcolor!15}\hrule}
    \vspace{1ex}

    \item[\faCode] \textbf{Agent 对话日志数据分析与评测体系} | Python, Pandas, NumPy, scikit-learn

    多 Agent 系统（Pi / Claude / ChatGPT / DeepSeek）的对话日志包含大量非结构化文本，行为表现缺乏量化评测框架。
    \begin{itemize}
        \item 搭建对话日志数据处理管线，解决多源日志交错解析、噪声清洗等数据工程问题
        \item 基于统计分析对多模型进行对比测试，建立行为对比评测框架，从情绪一致性、风格切换阈值、跨会话影响等维度产出量化洞察
        \item 提出三层分级记忆模型与动态摘要缓存的数据方案，分析脚本可扩展用于任意 Agent 日志的行为模式提取
    \end{itemize}
\end{itemize}

%======================================================================
\sectionTitle{工作经历}{\faBriefcase}
%======================================================================
\begin{experiences}
  \experience%
    [2025.12]%
    {2026.05}%
    {广西一三数据科技有限公司 \quad AI训练师}%
    [\begin{itemize}
      \item 主导多模态 AIGC 数据评估体系建设，制定覆盖 15 类场景的评估标准与质量管控流程，将团队交付合格率提升至 98\%
      \item 负责模型输出质量分析，累计处理并产出高质量偏好数据 100 万+条，支撑模型训练与能力迭代
      \item 主导自动驾驶 3D 感知数据的标注策略制定与质量管控，处理叠帧重影和移动物体轨迹连续性问题，累计处理 100 万+ 目标物
    \end{itemize}]
\end{experiences}

%======================================================================
\sectionTitle{教育背景}{\faGraduationCap}
%======================================================================
\begin{educations}
  \education%
    {2020.09}%
    [2023.06]%
    {桂林电子科技大学}%
    {数学与计算科学学院}%
    {数学}%
    {硕士}
  \separator{0.2ex}
  \education%
    {2016.09}%
    [2020.06]%
    {西安邮电大学}%
    {理学院}%
    {信息与计算科学}%
    {学士}
\end{educations}

\end{document}
